\chapter{Analise dos benchmarks de operacoes v3}
\section{Contrato experimental}
Foram analisados 2340 registros, com 3 repeticoes e 5 seeds de dados.
As comparacoes sao pareadas por entrada, seed, perfil, operacao, repeticao e grupo de baseline.
\section{Desempenho}
Classificacoes sustentadas: 0; mistas: 10; contraditas: 0.
Uma alegacao de hardware exige speedup mediano de pelo menos 1,05, intervalo de 95\% acima de 1 e kernel confirmado.
\section{Hipoteses}
\begin{itemize}
\item \texttt{H1\_H5}: \texttt{int8\_vs\_pruning\_10} -- \texttt{sustentada}.
\item \texttt{H1\_H5}: \texttt{int8\_vs\_pruning\_25} -- \texttt{sustentada}.
\item \texttt{H1\_H5}: \texttt{int8\_vs\_pruning\_50} -- \texttt{sustentada}.
\item \texttt{H1\_H5}: \texttt{int8\_vs\_pruning\_75} -- \texttt{sustentada}.
\item \texttt{H6}: \texttt{mse\_vs\_sparsity} -- \texttt{sustentada}.
\end{itemize}
\section{Limites}
Resultados numericos sem kernel confirmado nao sustentam ganho fisico. Os intervalos bootstrap e todos os casos individuais permanecem nos CSVs.
